\documentclass[11pt, letterpaper]{article}

% ── packages ──────────────────────────────────────────────────────────────────
\usepackage[margin=1in]{geometry}
\usepackage[T1]{fontenc}
\usepackage{lmodern}
\usepackage{microtype}
\usepackage{xcolor}
\usepackage{titlesec}
\usepackage{enumitem}
\usepackage{booktabs}
\usepackage{array}
\usepackage{longtable}
\usepackage{multirow}
\usepackage{tabularx}
\usepackage{hyperref}
\usepackage{parskip}
\usepackage{mdframed}
\usepackage{fancyhdr}
\usepackage{graphicx}
\usepackage{amsmath}

% ── colors ────────────────────────────────────────────────────────────────────
\definecolor{accent}{HTML}{1f6f78}
\definecolor{accent2}{HTML}{e07a5f}
\definecolor{lightgray}{HTML}{f5f5f5}
\definecolor{darkgray}{HTML}{2d2d2d}
\definecolor{mutedgray}{HTML}{5c6672}

% ── hyperref setup ────────────────────────────────────────────────────────────
\hypersetup{
    colorlinks = true,
    linkcolor  = accent,
    urlcolor   = accent,
    citecolor  = accent,
    pdftitle   = {Music Replayability Analysis — Project Report},
    pdfauthor  = {Chinmay Govind, Rohith Yelisetty},
}

% ── section formatting ────────────────────────────────────────────────────────
\titleformat{\section}
    {\large\bfseries\color{accent}}
    {\thesection.}{0.5em}{}
    [\vspace{-0.3em}\color{accent}\rule{\linewidth}{0.6pt}\vspace{0.2em}]

\titleformat{\subsection}
    {\normalsize\bfseries\color{darkgray}}
    {\thesubsection.}{0.5em}{}

\titleformat{\subsubsection}
    {\normalsize\itshape\color{mutedgray}}
    {\thesubsubsection.}{0.5em}{}

% ── page style ────────────────────────────────────────────────────────────────
\pagestyle{fancy}
\fancyhf{}
\renewcommand{\headrulewidth}{0.4pt}
\fancyhead[L]{\small\color{mutedgray}Music Replayability Analysis}
\fancyhead[R]{\small\color{mutedgray}CIS 2450 — Big Data Analytics}
\fancyfoot[C]{\small\thepage}

% ── highlight box ─────────────────────────────────────────────────────────────
\newmdenv[
    backgroundcolor = lightgray,
    linecolor       = accent,
    linewidth       = 1.5pt,
    leftline        = true,
    rightline       = false,
    topline         = false,
    bottomline      = false,
    innerleftmargin = 10pt,
    innerrightmargin= 8pt,
    innertopmargin  = 6pt,
    innerbottommargin= 6pt,
]{infobox}

% ── monospace file-tree helpers ───────────────────────────────────────────────
\newcommand{\ftree}[1]{\texttt{\small #1}}
\newcommand{\fdesc}[1]{\small\color{mutedgray} — #1}

% ── document ──────────────────────────────────────────────────────────────────
\begin{document}

% ── title page ────────────────────────────────────────────────────────────────
\begin{center}
    {\LARGE\bfseries\color{accent} Music Replayability Analysis}\\[0.5em]
    {\large Predicting Song Replay Behavior from Metadata, Audio Features, and Lyrics}\\[1em]
    {\normalsize CIS 2450: Big Data Analytics \quad|\quad Spring 2026 Final Project}\\[0.3em]
    {\normalsize Chinmay Govind\ (\href{mailto:cgovind@seas.upenn.edu}{cgovind@seas.upenn.edu}) \quad\textbullet\quad Rohith Yelisetty\ (\href{mailto:rohithy@seas.upenn.edu}{rohithy@seas.upenn.edu})}\\[0.5em]
    {\small \color{mutedgray}\today}
\end{center}

\vspace{0.5em}
\color{accent}\rule{\linewidth}{1pt}\color{black}
\vspace{1em}

% ─────────────────────────────────────────────────────────────────────────────
\section{Project Summary}
% ─────────────────────────────────────────────────────────────────────────────

This project investigates the factors that make songs highly replayable, using
data collected from three public music APIs and a large-scale lyrics dataset.
The target variable, \textbf{repeat\_listens}, is derived as the difference between
total listen count and total unique listener count from ListenBrainz,
representing listens beyond a song's first play.

The pipeline spans the full data science lifecycle: API-driven data collection,
multi-source joining and feature engineering, exploratory data analysis,
regression and classification modeling, NLP-based lyrics analysis, and an
interactive Streamlit dashboard. The best regression model (tuned Gradient
Boosting) achieved a holdout $R^2 = 0.3118$, showing that structured metadata
alone carries meaningful signal for predicting replayability.

\begin{infobox}
\textbf{Key Findings:}
\begin{itemize}[noitemsep, topsep=2pt]
    \item \textbf{55,579} unique tracks after cleaning, across \textbf{20 genres} and \textbf{14,555} artists
    \item \textbf{3 API sources:} MusicBrainz (metadata) + ListenBrainz (popularity) + AcousticBrainz (audio)
    \item \textbf{1 offline source:} Kaggle 5M Song Lyrics Dataset (lyrics text analysis)
    \item \textbf{20,759} songs matched to lyrics for NLP analysis
    \item Best model: Gradient Boosting (tuned) — $R^2 = 0.3118$, RMSE $= 2.37$
\end{itemize}
\end{infobox}

% ─────────────────────────────────────────────────────────────────────────────
\section{Codebase Overview}
% ─────────────────────────────────────────────────────────────────────────────

The codebase is organized as a set of standalone Python modules, each
responsible for one stage of the pipeline. All shared paths and constants live
in \texttt{config.py} so no path is hardcoded elsewhere.

\subsection{File Tree}

\begin{mdframed}[backgroundcolor=lightgray, linewidth=0pt,
                 innerleftmargin=12pt, innertopmargin=8pt, innerbottommargin=8pt]
\begin{small}\ttfamily
CIS2450FinalProject/\\
\quad config.py \fdesc{shared constants, paths, feature lists, targets}\\
\quad data\_collection.py \fdesc{API fetching: MusicBrainz, ListenBrainz, AcousticBrainz}\\
\quad data\_processing.py \fdesc{join, feature engineering, DuckDB SQL analytics}\\
\quad eda.py \fdesc{8 EDA plots + Kruskal-Wallis and Spearman tests}\\
\quad models.py \fdesc{7 regression models, PCA, hyperparameter tuning, imbalance demo}\\
\quad lyrics\_analysis.py \fdesc{lyrics loading, NLP features, LDA, 10 plots}\\
\quad dashboard.py \fdesc{Streamlit dashboard — 6 tabs, live prediction console}\\
\quad requirements.txt \fdesc{pinned Python dependencies}\\[0.4em]
\quad cache/ \fdesc{pickle caches for API responses and lyrics}\\
\quad\quad musicbrainz\_cache.pkl\\
\quad\quad listenbrainz\_cache.pkl\\
\quad\quad acousticbrainz\_cache.pkl\\
\quad\quad kaggle\_lyrics\_cache.pkl \fdesc{built on first run of lyrics\_analysis.py}\\[0.4em]
\quad outputs/ \fdesc{all generated artifacts}\\
\quad\quad music\_dataset\_processed.csv \fdesc{final 55 k-row processed dataset}\\
\quad\quad eda1\_*.png\ \ldots\ eda8\_*.png \fdesc{EDA visualizations}\\
\quad\quad lyrics1\_*.png\ \ldots\ lyrics11\_*.png \fdesc{lyrics NLP + model visualizations}\\
\quad\quad model\_results.csv \fdesc{CV + holdout metrics for all 7 models}\\
\quad\quad best\_gbm\_model.joblib \fdesc{serialized production model}\\
\quad\quad lyrics\_features.csv \fdesc{per-song NLP features}\\
\quad\quad eda\_summary.md \quad model\_summary.md \quad lyrics\_summary.md\\[0.4em]
\quad lyrics/ \fdesc{place the Kaggle CSV here before running lyrics\_analysis.py}\\
\quad\quad ds2.csv \fdesc{5M Song Lyrics Dataset (9.2 GB, not committed to git)}\\[0.4em]
\quad report/\\
\quad\quad report.tex \fdesc{this document}
\end{small}
\end{mdframed}

\subsection{Module Descriptions}

\begin{longtable}{>{\ttfamily\small}p{4.2cm} p{10.3cm}}
\toprule
\normalfont\textbf{File} & \textbf{Responsibility} \\
\midrule
config.py &
    Single source of truth for all path constants (\texttt{PROCESSED\_CSV},
    \texttt{CACHE\_DIR}, etc.), feature lists (\texttt{NUMERIC\_FEATURES},
    \texttt{AUDIO\_NUMERIC\_FEATURES}), and modeling targets
    (\texttt{TARGET = "log\_repeat\_listens"}). \\[4pt]

data\_collection.py &
    Fetches recordings from MusicBrainz by genre tag (up to 5\,000 per genre),
    enriches them with ListenBrainz popularity metrics in batches of 500, and
    optionally adds AcousticBrainz audio features via a multi-threaded fetcher.
    All results are pickled for fast re-runs. \\[4pt]

data\_processing.py &
    Deduplicates MusicBrainz rows, performs a 3-way join (MusicBrainz +
    ListenBrainz + AcousticBrainz) in Polars, engineers new features, winsorizes
    outliers, label-encodes categoricals, and runs 8 analytical SQL queries
    in DuckDB. Writes the final processed CSV and a JSON data summary. \\[4pt]

eda.py &
    Generates 8 diagnostic plots (target distribution, genre ranking, decade
    trend, tempo/danceability, correlation heatmap, duration buckets,
    missingness, outliers by genre) and runs Kruskal-Wallis and Spearman
    statistical tests. \\[4pt]

models.py &
    Trains 7 models (Linear Regression baseline, Ridge, Lasso, PCA+Ridge,
    Random Forest, Gradient Boosting, and a tuned GBM via RandomizedSearchCV),
    evaluates with 5-fold CV and holdout metrics, extracts feature importances,
    runs a classification imbalance demo, and serializes the best model. \\[4pt]

lyrics\_analysis.py &
    Streams the 9.2 GB Kaggle CSV in 50 k-row chunks to match all processed
    songs by normalized artist/title, caches results as a pickle, then computes
    VADER sentiment, type-token ratio, repetitiveness, Flesch readability,
    bigram frequencies, word length statistics, and rare-word ratios
    (via \texttt{wordfreq}). Fits an 8-topic LDA model and produces 10 plots. \\[4pt]

dashboard.py &
    Streamlit application with 6 tabs: Command Center (dataset overview),
    Genre Lab (interactive drill-down), Model Studio (model comparison and
    feature importance), Prediction Console (live scoring form), Lyrics Analysis
    (all NLP plots and tables), and Song Explorer (per-song autocomplete
    lookup with live Genius API lyric fetch and per-song NLP breakdown). \\

\bottomrule
\end{longtable}

% ─────────────────────────────────────────────────────────────────────────────
\section{Dataset Overview}
% ─────────────────────────────────────────────────────────────────────────────

\subsection{Source 1 — MusicBrainz (Recording Metadata)}

MusicBrainz is a community-maintained music encyclopedia with a free,
open REST API (\texttt{musicbrainz.org/ws/2/recording}). For each of 20
genre tags the project issues paginated search queries, collecting up to
5\,000 recordings per genre. Each recording provides: MBID (a permanent
unique identifier), title, duration, release year, release type (Album /
Single / EP), and artist metadata (name, type, country, career start year).

Songs that appear under multiple genre searches are deduplicated on MBID;
all their genre tags are retained and collapsed into a pipe-separated string.

\subsection{Source 2 — ListenBrainz (Popularity Signals)}

ListenBrainz tracks listening activity for millions of users.
The popularity endpoint (\texttt{api.listenbrainz.org/1/popularity/recording})
accepts batches of 500 MBIDs and returns \texttt{total\_listen\_count} and
\texttt{total\_user\_count} for each. The modeling target is defined as:

\[
\texttt{repeat\_listens} = \texttt{total\_listen\_count} - \texttt{total\_user\_count}
\]

representing listens beyond one per unique listener — i.e.\ genuine replay
behavior. Songs with no ListenBrainz record are excluded.

\subsection{Source 3 — AcousticBrainz (Audio Features) — Optional Enrichment}

AcousticBrainz stored Essentia-derived audio descriptors.
The project fetches low-level features per MBID (with 6 concurrent threads)
extracting: tempo (BPM), danceability, musical key and scale, average loudness,
and dynamic complexity. Coverage is 42.81\% of the processed dataset;
the modeling pipeline therefore treats audio features as optional enrichment
rather than mandatory inputs.

\subsection{Source 4 — Kaggle 5M Song Lyrics Dataset}

The Kaggle dataset (\textit{nikhilnayak123/5-million-song-lyrics-dataset})
is a 9.2 GB CSV containing lyrics for approximately 5 million songs.
The lyrics analysis module streams it in 50 k-row chunks and matches each row
to the processed dataset by normalizing both artist name and title
(lowercase, strip punctuation). Matched lyrics are cached as a pickle so
the 9.2 GB file is scanned only once. \textbf{20,759 songs} were matched.

\subsection{Final Dataset Statistics}

\begin{center}
\begin{tabular}{lr}
\toprule
\textbf{Attribute} & \textbf{Value} \\
\midrule
Total tracks (after cleaning) & 55,579 \\
Total columns & 39 \\
Unique artists & 14,555 \\
Genres & 20 \\
Median repeat listens & 344 \\
Median log repeat listens & 5.84 \\
AcousticBrainz coverage & 42.81\% \\
Songs matched to lyrics & 20,759 \\
\bottomrule
\end{tabular}
\end{center}

% ─────────────────────────────────────────────────────────────────────────────
\section{Data Processing Overview}
% ─────────────────────────────────────────────────────────────────────────────

\subsection{Joining and Deduplication}

The three API sources are joined in Polars on MBID. MusicBrainz rows are
first deduplicated: songs fetched under multiple genres have their genre tags
aggregated, and the primary genre is taken as the alphabetically first tag.
The \texttt{genre\_match\_count} column records how many genre searches each
song appeared in (used as a feature).

\subsection{Feature Engineering}

The following features are derived during processing:

\begin{itemize}[noitemsep]
    \item \textbf{duration\_sec} — converted from milliseconds, clipped to [0, 900]
    \item \textbf{release\_decade} — floored release year (e.g.\ 1990, 2000)
    \item \textbf{track\_age} — \texttt{current\_year} $-$ \texttt{release\_year}
    \item \textbf{artist\_career\_age} — release year $-$ artist debut year
    \item \textbf{log\_repeat\_listens} — $\ln(\texttt{repeat\_listens} + 1)$, the regression target
    \item \textbf{has\_audio\_features} — binary flag for AcousticBrainz coverage
    \item \textbf{audio\_feature\_missing\_count} — count of null audio fields (0 or 4)
    \item \textbf{tempo\_x\_dance} — interaction term: tempo $\times$ danceability
    \item \textbf{career\_x\_duration\_min} — interaction term: career age $\times$ duration in minutes
    \item \textbf{is\_high\_replay} — classification target (top-25\% repeat listens = 1)
\end{itemize}

\subsection{Outlier Handling}

Continuous features (tempo, danceability, loudness, dynamic complexity,
duration\_sec) are winsorized at the 1st and 99th percentiles to limit the
influence of extreme values without discarding records.

\subsection{Categorical Encoding}

Six categorical columns (genre, release\_type, artist\_type, artist\_country,
key, key\_scale) are label-encoded into integer columns suffixed \texttt{\_enc}.
Original string columns are retained for display and filtering in the dashboard.

\subsection{Null Handling}

Missing values in numeric features are handled inside scikit-learn pipelines
using \texttt{SimpleImputer} (median strategy), ensuring no statistics are
computed on test data and no leakage occurs.

\subsection{SQL Analytics with DuckDB}

After building the Polars frame, \texttt{data\_processing.py} registers it
as an in-memory DuckDB table and runs 8 SQL queries covering:
genre-level summary statistics, decade trends, top artists, country breakdown,
duration buckets, genre overlap, and AcousticBrainz coverage. Results are
exported as individual CSVs and consumed by the dashboard.

\subsection{Entity Linking}

The lyrics module performs a form of entity linking by normalizing
artist and title strings on both sides (lowercase, strip punctuation,
collapse whitespace) and performing an exact lookup against the Kaggle dataset.
This links the internal MBID-keyed dataset to an external source that has no
MBID column — demonstrating cross-dataset record reconciliation.

% ─────────────────────────────────────────────────────────────────────────────
\section{Models Overview}
% ─────────────────────────────────────────────────────────────────────────────

\subsection{Experimental Setup}

\begin{itemize}[noitemsep]
    \item \textbf{Train/test split:} 80\,/\,20 stratified on the classification target
    \item \textbf{Cross-validation:} 5-fold on the training split only
    \item \textbf{Imputation:} median imputation inside each fold's pipeline
    \item \textbf{Metrics:} RMSE, MAE, and $R^2$ for regression; accuracy, F1, ROC-AUC for classification
    \item \textbf{Feature set:} metadata-first; AcousticBrainz features included where available
\end{itemize}

\subsection{Models Trained}

\begin{enumerate}[noitemsep]
    \item \textbf{Linear Regression} (baseline) — ordinary least squares, no regularization
    \item \textbf{Ridge Regression} — L2 regularization, $\alpha = 1.0$
    \item \textbf{Lasso Regression} — L1 regularization, $\alpha = 0.1$
    \item \textbf{PCA + Ridge} — 16 principal components (95.2\% explained variance) then Ridge; explores dimensionality reduction
    \item \textbf{Random Forest} — ensemble of 200 decision trees
    \item \textbf{Gradient Boosting} — default scikit-learn GBM
    \item \textbf{Gradient Boosting (tuned)} — \texttt{RandomizedSearchCV} over 5 parameters: learning rate, number of estimators, max depth, subsampling rate, min samples per leaf
\end{enumerate}

\subsection{Difficulty Concepts Applied}

\begin{itemize}[noitemsep]
    \item \textbf{Feature importance} — extracted from tree models (Random Forest and both GBM variants); top signals include genre (classical, indie), duration, track age, and release year
    \item \textbf{PCA for feature selection} — 16 components retained from the full feature matrix, explaining 95.2\% of variance; compared against direct regression
    \item \textbf{Hyperparameter tuning} — \texttt{RandomizedSearchCV} on the GBM with 20 candidates and 5-fold CV; best params: $n\_\text{estimators}=500$, $\text{learning\_rate}=0.05$, $\text{max\_depth}=5$, $\text{subsample}=0.7$
    \item \textbf{Ensemble models} — Random Forest and Gradient Boosting both outperform all linear baselines
    \item \textbf{Imbalance handling} — logistic regression trained with and without \texttt{class\_weight=\textquotesingle{}balanced\textquotesingle{}}; balanced weights improve F1 from 0.327 to 0.506 at the cost of accuracy (0.766 $\to$ 0.691) while ROC-AUC is unchanged (0.737)
    \item \textbf{Feature engineering} — interaction terms (\texttt{tempo\_x\_dance}, \texttt{career\_x\_duration\_min}) and derived temporal features (\texttt{track\_age}, \texttt{artist\_career\_age}, \texttt{release\_decade})
\end{itemize}

% ─────────────────────────────────────────────────────────────────────────────
\section{Results Overview}
% ─────────────────────────────────────────────────────────────────────────────

\subsection{Regression Model Comparison}

\begin{center}
\begin{tabular}{lrrr}
\toprule
\textbf{Model} & \textbf{CV $R^2$} & \textbf{Holdout $R^2$} & \textbf{Holdout RMSE} \\
\midrule
Linear Regression (baseline) & 0.247 & 0.232 & 2.505 \\
Ridge                         & 0.247 & 0.232 & 2.505 \\
Lasso                         & 0.247 & 0.232 & 2.506 \\
PCA + Ridge                   & 0.185 & 0.168 & 2.607 \\
Random Forest                 & 0.311 & 0.301 & 2.391 \\
Gradient Boosting             & 0.307 & 0.294 & 2.403 \\
\textbf{Gradient Boosting (tuned)} & \textbf{0.321} & \textbf{0.312} & \textbf{2.372} \\
\bottomrule
\end{tabular}
\end{center}

\vspace{0.5em}

The tuned Gradient Boosting model is the clear winner, outperforming all
linear baselines by approximately 8 $R^2$ points. Ensemble methods capture
the non-linear interaction between genre, track age, duration, and
replay behavior that linear models cannot represent. PCA compression
slightly hurts performance here, suggesting the original feature space
does not contain severe multicollinearity that needs compressing.

\subsection{Top Feature Importances (Tuned GBM)}

\begin{center}
\begin{tabular}{lr}
\toprule
\textbf{Feature} & \textbf{Importance} \\
\midrule
genre: classical      & 0.197 \\
genre: indie          & 0.133 \\
duration\_sec         & 0.131 \\
track\_age            & 0.095 \\
release\_year         & 0.089 \\
genre: r\&b           & 0.062 \\
audio\_feature\_missing\_count & 0.049 \\
genre: punk           & 0.044 \\
has\_audio\_features  & 0.039 \\
genre: metal          & 0.023 \\
\bottomrule
\end{tabular}
\end{center}

Genre is the dominant predictor, particularly classical and indie — genres
with concentrated, committed listener bases. Duration and track age together
suggest that well-aged mid-length songs accumulate the most replays over time.

\subsection{EDA Findings}

\begin{itemize}[noitemsep]
    \item \textbf{Target distribution:} \texttt{repeat\_listens} is heavily right-skewed; the log transform stabilizes variance and is essential for regression.
    \item \textbf{Genre effect:} Kruskal-Wallis test: $H = 6395$, $p \approx 0$ — replay distributions differ significantly across genres.
    \item \textbf{Duration effect:} Spearman $\rho = 0.042$, $p < 10^{-21}$ — mid-length tracks (3–5 min) replay best; very short and very long tracks lag.
    \item \textbf{Decade trend:} Replayability has shifted over release eras, motivating \texttt{release\_decade} and \texttt{track\_age} features.
    \item \textbf{Audio features:} Missingness is concentrated in AcousticBrainz fields (57\%), confirming the metadata-first modeling strategy.
\end{itemize}

\subsection{Lyrics Analysis Findings}

\textbf{23,641 songs} were matched to lyrics from the Kaggle dataset (using exact + fuzzy artist/title matching via \texttt{rapidfuzz}).
Key aggregate findings:

\begin{itemize}[noitemsep]
    \item Average VADER compound sentiment: $+0.190$ (mildly positive overall)
    \item Average vocabulary richness (TTR): 0.539
    \item Average repetitiveness (duplicate-line fraction): 0.231
    \item Sentiment has the strongest lyric-level correlation with replay: $r = -0.088$ (slightly more negative songs replay more)
    \item \textbf{Top words in high-replay songs:} dead, blood, watch, fuck, black, shit, fall, war, fire, inside
    \item \textbf{Top bigrams in high-replay songs:} ``lê lê'', ``close eyes'', ``take back'', ``never get'', ``run away''
    \item \textbf{LDA topics:} 8 coherent themes — multilingual (``que de ich und''), conversational (``don't got know''), romantic (``love baby girl''), hip-hop (``shit fuck nigga''), spiritual (``god lord soul''), rootsy (``good man home town''), dance/rock (``night dance sing rock''), temporal (``never time away day'')
\end{itemize}

\subsubsection*{Lyrics Prediction Model}

To quantify how much lyrical content predicts replay behavior, we trained regression models on three feature sets: \textbf{Lyrics Only} (9 NLP features + 8 LDA topic weights), \textbf{Metadata Only} (numeric + categorical metadata on the same matched subset for a fair baseline), and \textbf{Lyrics + Metadata} (combined). Both Ridge regression and Gradient Boosting (GBM) were evaluated with 5-fold CV and a shared 80/20 holdout split.

\begin{center}
\begin{tabular}{llccc}
\toprule
Feature Set & Model & CV $R^2$ & Holdout $R^2$ & Holdout RMSE \\
\midrule
Lyrics Only        & Ridge & $0.024 \pm 0.002$ & 0.033 & 2.502 \\
Lyrics Only        & GBM   & $0.034 \pm 0.003$ & 0.044 & 2.488 \\
Metadata Only      & Ridge & $0.020 \pm 0.003$ & 0.022 & 2.516 \\
Metadata Only      & GBM   & $0.193 \pm 0.008$ & 0.205 & 2.270 \\
Lyrics + Metadata  & Ridge & $0.041 \pm 0.002$ & 0.053 & 2.476 \\
Lyrics + Metadata  & GBM   & $0.191 \pm 0.008$ & 0.208 & 2.265 \\
\bottomrule
\end{tabular}
\end{center}

Key takeaways: lyrics features alone explain only $\sim$3–4\% of variance, while metadata GBM reaches $\sim$20\%. Combining both yields a marginal improvement ($R^2 = 0.208$), meaning lyrical signals provide modest incremental signal beyond metadata but are not the primary driver of replayability in this dataset.

% ─────────────────────────────────────────────────────────────────────────────
\section{Dashboard Overview}
% ─────────────────────────────────────────────────────────────────────────────

The dashboard is built with Streamlit (\texttt{streamlit run dashboard.py})
and organized into six tabs, each covering a distinct analytical workspace.

\begin{longtable}{>{\bfseries\small}p{3.2cm} p{11.3cm}}
\toprule
Tab & Description \\
\midrule
Command Center &
    Top-line dataset health: track count, artist count, audio coverage, best
    model card; genre distribution bar chart; decade replayability trend;
    top-25 artist leaderboard; AcousticBrainz coverage by genre;
    collapsible EDA and modeling narrative summaries. \\[4pt]

Genre Lab &
    Interactive controls (genre multiselect, decade range slider, scatter
    sample size) that filter the processed dataset in real time. Shows
    replayability trend lines per genre, a box-plot spread comparison,
    a duration scatter, and a filtered artist leaderboard. \\[4pt]

Model Studio &
    Cross-validation vs.\ holdout $R^2$ grouped bar chart; holdout RMSE/MAE
    comparison; interactive feature importance view (selectable model);
    tuned GBM parameters JSON; PCA summary; actual vs.\ predicted scatter;
    imbalance handling results table. \\[4pt]

Prediction Console &
    Live scoring form that accepts a full track profile (genre, release type,
    artist type, country, key, duration, year, career age, optional audio
    features) and returns the predicted log repeat listens, estimated raw
    replays, and dataset percentile. \\[4pt]

Lyrics Analysis &
    Four stat cards (tracks analyzed, avg sentiment, vocabulary richness,
    repetitiveness); all 11 NLP/model plots in a 2-column grid (including
    lyrics prediction model comparison); model results table; top-30 word
    frequency comparison table; collapsible full summary report. \\[4pt]

Song Explorer &
    Typeahead selectbox over up to 10\,000 songs (top by replay score).
    Selecting a track shows its metadata cards, checks the local lyrics cache,
    falls back to the Genius API if not cached, then renders a per-song
    NLP analysis: sentiment breakdown chart, top-15 words bar chart,
    vocabulary/repetitiveness/readability metrics, and a full lyrics expander. \\

\bottomrule
\end{longtable}

\subsection{Course Topics Demonstrated in the Dashboard}

\begin{itemize}[noitemsep]
    \item \textbf{Polars} — all data loading and filtering in the dashboard uses Polars DataFrames
    \item \textbf{SQL (DuckDB)} — 8 analytical queries power the Command Center charts
    \item \textbf{Text representations \& NLP} — VADER sentiment, type-token ratio, LDA topic model, bigram frequencies, \texttt{wordfreq} rarity scoring
    \item \textbf{Record linking} — MusicBrainz MBID links metadata to popularity; normalized artist/title matching links to the Kaggle lyrics corpus
    \item \textbf{Unsupervised learning} — 8-topic LDA surfaced in the Lyrics Analysis tab
    \item \textbf{Supervised learning} — 7 regression models and a classification demo in Model Studio
    \item \textbf{Joins} — 3-way equi-join on MBID in the processing pipeline
    \item \textbf{Feature engineering} — interaction terms, temporal features, audio enrichment flags

\end{itemize}

% ─────────────────────────────────────────────────────────────────────────────
\section{Conclusion}
% ─────────────────────────────────────────────────────────────────────────────

This project demonstrates that song replayability can be meaningfully
predicted from freely available metadata and listening signals, without
requiring proprietary audio analysis. The tuned Gradient Boosting model
achieves $R^2 = 0.312$ on the holdout set — a substantial improvement over
the linear baseline ($R^2 = 0.232$) — primarily driven by genre,
duration, and track age.

\textbf{Key takeaways:}
\begin{itemize}[noitemsep]
    \item Genre is the single strongest predictor of replay behavior, reflecting audience
          loyalty patterns rather than intrinsic audio properties.
    \item Duration has a non-linear effect: 3–5 minute tracks replay best, consistent
          with the industry-standard ``radio edit'' convention.
    \item Sentiment has a small but consistent negative correlation with replay
          ($r = -0.089$): darker lyrical themes are modestly associated with higher replays.
    \item Vocabulary richness (TTR) and repetitiveness show weaker correlations,
          suggesting lyrical complexity plays a secondary role to thematic content.
    \item The LDA topic model reveals that high-replay songs cluster around nocturnal,
          relationship, and intensity-driven themes, while spiritual and non-English
          topics dominate lower-replay clusters.
\end{itemize}

\textbf{Limitations and future work:}
\begin{itemize}[noitemsep]
    \item AcousticBrainz is a defunct service (shut down 2022); 57\% of tracks have
          no audio features, capping the model's ability to exploit tempo and danceability.
    \item Artist/title normalization achieves $\sim$37\% lyrics match rate — fuzzy
          matching could recover additional songs.
    \item The regression target is a proxy: ListenBrainz captures only users who
          actively scrobble, biasing toward niche and enthusiast audiences.
    \item Future extensions could incorporate embedding-based lyric representations
          (e.g.\ sentence transformers) or time-series replay trajectories to capture
          temporal popularity decay.
\end{itemize}

\vspace{1em}
\color{accent}\rule{\linewidth}{0.6pt}\color{black}
\begin{center}
\small\color{mutedgray}
Full codebase at \url{https://github.com/rohithyelisetty/CIS2450FinalProject}
\quad|\quad
Run order: \texttt{data\_collection.py} $\to$ \texttt{data\_processing.py} $\to$
\texttt{eda.py} $\to$ \texttt{models.py} $\to$ \texttt{lyrics\_analysis.py} $\to$
\texttt{streamlit run dashboard.py}
\end{center}

\end{document}
